\originalchapter{8.06 作业9：交换相关与有效自旋}
来源为Spring 2018 Assignment 9 PRACTICE, 原件发布行写Thursday May 15 2018（该日期实际是星期二）, 按日期记录而不静默改原文星期. 这是练习集, 第1、2题中文转述及\textbf{AI独立解答}. 第3题署Inspired by Cohen--Tannoudji, 引用不自动等于复制; 本轮具体表达尚核, 未计入解题完成数, 不整份禁入.
\originalsection{第1题：交换“力”是相关}
\begin{exercise}
(a)正交单粒子态$\alpha,\beta$, 对乘积态$D$、归一对称态$S$、反对称态$A$, 求$\langle(x_1-x_2)^2\rangle$之差并解释符号. (b)两spin-$s$的$m_a\ne m_b$态也作三组合, 求$\langle(S_1-S_2)^2\rangle$及排序.
\end{exercise}
\begin{solution}
设$X=\bra\alpha x\ket\beta$. 单粒子平方项的交换交叉都含$\braket\alpha\beta=0$, 只有$-2x_1x_2$产生交叉$-2|X|^2$. 故$D_D=\langle x^2\rangle_\alpha+\langle x^2\rangle_\beta-2\langle x\rangle_\alpha\langle x\rangle_\beta$, $D_S-D_D=-2|X|^2$, $D_A-D_D=+2|X|^2$. 对称增强靠近相关, 反对称抑制; 没有新势能力作用, 是交换对称改变波幅相干.

自旋$D_D=2\hbar^2[s(s+1)-m_am_b]$. 交换交叉的$S_{1z}S_{2z}$为0; $S_1\cdot S_2=S_{1z}S_{2z}+(S_{1+}S_{2-}+S_{1-}S_{2+})/2$, 只有$|m_a-m_b|=1$有非零交叉. 若相邻, 记$m_< =\min(m_a,m_b)$,$F=s(s+1)-m_<(m_<+1)>0$, 则$D_S=D_D-\hbar^2F$, $D_A=D_D+\hbar^2F$. 非相邻三者相等. 所以总有$D_S\le D_D\le D_A$, 等号条件不忽略; 类似位置相关却受升降选择规则限制.
\end{solution}
\originalsection{第2题：空间Coulomb作用如何表现为自旋作用}
\begin{exercise}
两电子占正交等能轨道$\psi_1,\psi_2$, 总态须反对称. (a)(b)空间反对称、对称分别配什么自旋态, 求$S^2,S_z$读数. (c)加$CS_1\cdot S_2$求能量. (d)改只有$e^2/r_{12}$Coulomb排斥, 以积分给一阶两组能差. (e)用$H_e=-JS_1\cdot S_2$表示相同能差并判断符号.
\end{exercise}
\begin{solution}
空间$\Psi_\pm=(\psi_1(r_1)\psi_2(r_2)\pm\psi_2(r_1)\psi_1(r_2))/\sqrt2$. $\Psi_-$须配自旋对称triplet:$\uparrow\uparrow,(\uparrow\downarrow+\downarrow\uparrow)/\sqrt2,\downarrow\downarrow$, $S^2=2\hbar^2$, $S_z=\hbar,0,-\hbar$. $\Psi_+$须配singlet$(\uparrow\downarrow-\downarrow\uparrow)/\sqrt2$, 两读数0. 设轨道总能$2E_0$, $S_1\cdot S_2=(S^2-3\hbar^2/2)/2$, 故triplet能$2E_0+C\hbar^2/4$, singlet$2E_0-3C\hbar^2/4$, 自旋耦合本身不混上述本征态.

Coulomb直接积分$D=\iint |\psi_1(r_1)|^2|\psi_2(r_2)|^2e^2/r_{12}\dd^3r_1\dd^3r_2$, 交换积分$K=\iint\psi_1^*(r_1)\psi_2^*(r_2)\psi_2(r_1)\psi_1(r_2)e^2/r_{12}\dd^3r_1\dd^3r_2$. 变换积分变量核$K$实, 两组能量$2E_0+D\mp K$（triplet取减, singlet取加）. 定义$\epsilon=E_S-E_T=2K$, 则$H_e$两组差$J\hbar^2$, $J=\epsilon/\hbar^2=2K/\hbar^2$. 另加共同能量常数才能复现绝对能, 不是只用$H_e$丢去轨道平均值.

依题目直觉, 反对称空间使更分离、排斥较小, 预期$K\ge0,J\ge0$, 有效铁磁号. 这里还能以Coulomb核Fourier$4\pi/k^2$正性核实$K\ge0$（适当有限能轨道）, 但不能将$\langle(x_1-x_2)^2\rangle$增大单独视为任何势下能量下降的严格证明. 有效自旋项由Pauli约束与空间交换产生, 与显式磁矩直接自旋力$C$有不同来源.
\end{solution}
